\section{Alberi Binari (Lab 18)}

\subsection{Cos'è un Albero Binario?}
È una struttura gerarchica (come le cartelle del PC).
Ogni elemento si chiama **Nodo** e può avere massimo due figli:
\begin{enumerate}
    \item **Left** (Sinistro)
    \item **Right** (Destro)
\end{enumerate}

\subsubsection{Immaginiamolo visivamente}
\begin{center}
\begin{tikzpicture}[
    level distance=1.5cm,
    sibling distance=2.5cm,
    every node/.style={circle, draw, minimum size=0.8cm, fill=niceBlue!10, font=\bfseries},
    edge from parent/.style={draw, thick, -latex}
]

\node {10}
    child { node {5}
        child { node {2} }
        child { node {7} }
    }
    child { node {15}
        child { node {12} }
        child { node {20} }
    };
\end{tikzpicture}
\end{center}

\subsubsection{Rappresentazione in JS}
Non servono classi complesse, bastano oggetti.
\begin{lstlisting}[language=JavaScript]
let nodoS = { val: 5, left: null, right: null };
let nodoD = { val: 10, left: null, right: null };
let root  = { val: 7, left: nodoS, right: nodoD };
\end{lstlisting}

\subsection{Come visitare un Albero (Traversals)}
Dobbiamo passare per tutti i nodi. Ci sono vari modi.

\subsubsection{1. Visita in Profondità (DFS)}
\begin{concept}{Strategie DFS}
Andiamo giù il più possibile, poi torniamo indietro (ricorsione).
\begin{itemize}
    \item **Pre-Order:** IO $\to$ SX $\to$ DX
    \item **In-Order:** SX $\to$ IO $\to$ DX
    \item **Post-Order:** SX $\to$ DX $\to$ IO
\end{itemize}
\end{concept}

\begin{jsbox}
// Pre-Order: IO -> SX -> DX
function preOrder(nodo) {
    if (!nodo) return;
    console.log(nodo.val);
    preOrder(nodo.left);
    preOrder(nodo.right);
}

// In-Order: SX -> IO -> DX
function inOrder(nodo) {
    if (!nodo) return;
    inOrder(nodo.left);
    console.log(nodo.val);
    inOrder(nodo.right);
}

// Post-Order: SX -> DX -> IO
function postOrder(nodo) {
    if (!nodo) return;
    postOrder(nodo.left);
    postOrder(nodo.right);
    console.log(nodo.val);
}
\end{jsbox}

\subsubsection{2. Visita in Ampiezza (BFS)}
\begin{concept}{Strategia BFS (Queue)}
Visitiamo livello per livello (prima tutti i figli, poi tutti i nipoti).
Serve una **Coda (Queue)** per ricordarsi quali nodi visitare.
\end{concept}

\begin{jsbox}
function bfs(root) {
    let coda = [root]; // Inizio con la radice
    while (coda.length > 0) {
        let n = coda.shift(); // Prendo il primo
        console.log(n.val);
        
        if (n.left) coda.push(n.left);   // Aggiungo i figli in fondo
        if (n.right) coda.push(n.right);
    }
}
\end{jsbox}

\subsection{Alberi Binari di Ricerca (BST)}
Un albero speciale dove per ogni nodo:
\begin{itemize}
    \item Tutto ciò che è a \textbf{Sinistra} è \textbf{Minore}.
    \item Tutto ciò che è a \textbf{Destra} è \textbf{Maggiore}.
\end{itemize}

\subsubsection{Ricerca in un BST}
È velocissima (come la Binary Search negli array).
\begin{lstlisting}[language=JavaScript]
function search(nodo, x) {
    if (!nodo) return false;      // Non trovato
    if (nodo.val == x) return true; // Trovato!
    
    if (x < nodo.val) return search(nodo.left, x); // Cerca a sinistra
    return search(nodo.right, x); // Cerca a destra
}
\end{lstlisting}

\subsubsection{Inserimento in un BST}
Per inserire, scendiamo finché non troviamo un posto vuoto (`null`).
\begin{lstlisting}[language=JavaScript]
function insert(nodo, val) {
    if (!nodo) return { val: val, left: null, right: null };
    
    if (val < nodo.val) {
        nodo.left = insert(nodo.left, val);
    } else {
        nodo.right = insert(nodo.right, val);
    }
    return nodo;
}
\end{lstlisting}
